\documentclass[12pt]{article}
\usepackage[utf8]{inputenc}
\usepackage[T1]{fontenc}
\usepackage{hyperref}
\title{Yuanjiang-Mojiang Yi Language Resources for K-12 Curriculum Development}
\author{Compiled by Claude Sonnet 5.5 (claude-sonnet-5-5) for Andr\'as Kornai}
\date{}
\begin{document}
\maketitle

\section{Introduction}
The code yym (Glottolog yuan1242) is retired in ISO 639-3. Glottolog records that it was split into Southern Nisu [nsd] and Southwestern Nisu [nsv], because the two varieties are not mutually intelligible. The Glottolog entry gives the family as ``Bookkeeping'', a placeholder for retired codes. Searches found no resources under the name ``Yuanjiang-Mojiang Yi''. The resources below concern the Nisu varieties that replaced it. Yang (2009) places Southern Nisu in Honghe, Yuanyang, Jinping, Yuanjiang, Shiping and possibly Jianshui counties, and classes the Mojiang variety with Northern Nisu. This is a short report, and it should be redone separately for nsd and nsv.

\section{Primary Educational Resources}
No primers, textbooks or school programmes were found. A mission profile says an orthography was created but was accepted only in the area where it was produced, because speakers elsewhere use different vocabulary.

\section{Linguistic and Reference Materials}
\begin{itemize}
\item Glottolog entry for Yuanjiang-Mojiang Nisu (retired). \url{https://glottolog.org/resource/languoid/id/yuan1242}
\item Wikipedia, ``Nisu language''. \url{https://en.wikipedia.org/wiki/Nisu_language}
\item Wikipedia, ``Nisoish languages''. \url{https://en.wikipedia.org/wiki/Nisoish_languages}
\item Omniglot, Nisu writing: covers the traditional Yi script and the Latin alphabet. \url{https://omniglot.com/writing/nisu.htm}
\end{itemize}

\section{Religious and Community Resources}
\begin{itemize}
\item Asia Harvest people-group profile, Nisu (Yuanyang). \url{https://people-groups.asiaharvest.org/people-group-profiles/yunnan/nisu-yuanyang.pdf}
\item Global Recordings Network, Southern Nisu [nsd]. I did not verify what recordings it holds. \url{https://globalrecordings.net/en/language/nsd}
\end{itemize}

\section{Digital and Audio Resources}
Global Recordings Network also lists ``Yi: Lolo in Mojiang'', but under the code ysp (Lolopo, Southern), so it is probably a different language. \url{https://globalrecordings.net/en/language/6077}

\section{Curriculum Development Context}
Standardised Yi in China is usually taught from the Liangshan Yi script and schooling, and I found no Nisu-specific school materials. The next step would be to choose nsd or nsv and search Chinese-language sources for dictionaries and local Yi-language education in Honghe and Yuxi.
\end{document}